
\section{RLOP: Replication Learning of Option Pricing}

In this section, we propose a novel algorithm that prices a call/put
option via portfolio replication, but this method uses a forward view
which is fundamentally different from the QLBS approach. The idea
is simple: the agent manages a portfolio which yields a reward at
the terminal time based on how accurate the portfolio value is compared
to the option payoff. This naive idea has the problem that the reward
is zero for quite a long time until the maturity, which is usually
not good for shaping the agent's behavior. To deal with this downside,
we propose to group a few options as an ensemble so that the agent
gets a stream of feedbacks during each episode. To be specific, given
a (simulated or historical) path of the asset price $\left\{ S_{t}\right\} $,
maturity time $T$, and the payoff function $h$, the agent needs
to manages (at most) $T$ portfolios $\left\{ \Pi_{t}^{\left(i\right)}\right\} _{i=1}^{T}$
at the same time, where the $k$-th portfolio replicates the option
that terminals at time step $k$. In other words, for every $i\in\left[t+1,T\right]$,
the agent proposes the hedge position $u_{t}^{\left(i\right)}$ based
on the current time step $t$, the asset price $S_{t}$, and the balance
of the $i$-th portfolio $\Pi_{t}^{\left(i\right)}$. We hold the
belief that the agent is able to learn the hedging strategy step by
step from small $t$ to the terminal time $T$.

\subsection{MDP Formulation}

We now rigorously define the aforementioned problem as a MDP. Given
a maturity time $i$, the state space consists of tuples containing
the time step $t$, the asset price $S_{t}$, and the current portfolio
value $\Pi_{t}^{\left(i\right)}$. The action space is $\bR$ that
contains all possible hedging positions. The transitional probability
(density) function is defined as
\[
p\left(\left(t,S_{t},\Pi_{t}^{\left(i\right)}\right),u_{t}^{\left(i\right)}\to\left(t',S_{t'},\Pi_{t'}^{\left(i\right)}\right),R_{t+1}\right)=\begin{cases}
\delta_{t+1,t'}\delta_{\widetilde{\Pi}_{t+1}^{\left(i\right)},\Pi_{t'}^{\left(i\right)}}\rho\left(S_{t},S_{t'}\right) & t<i\\
0 & t=i
\end{cases}
\]
where 
\begin{itemize}
\item the only admissible state is the terminal state when $t>i$, marking
the end of this episode;
\item $\widetilde{\Pi}_{t+1}^{\left(i\right)}$ refers to the portfolio
value determined by the self-financing condition
\[
\widetilde{\Pi}_{t+1}^{\left(i\right)}=e^{r\Delta t}\left(\Pi_{t}-u_{t}^{\left(i\right)}S_{t}\right)+u_{t}^{\left(i\right)}S_{t+1}-\text{TC}\left(u_{t+1}^{\left(i\right)}-u_{t}^{\left(i\right)},S_{t+1}\right)
\]
(which is the same as Eqn. \ref{eq:qlbs-modified-self-financing});
\item $\rho$ characterizes the dynamics of the underlying asset, e.g. the
discrete version geometric brownian motion;
\item $R_{t+1}=0$ for $t+1<i$ and $R_{i}=H\left(h\left(S_{i}\right),\Pi_{i}^{\left(i\right)}\right)$
where $H$ is a given penalty function that measures how accurate
the portfolio value $\Pi_{i}^{\left(i\right)}$ mimics the option
payoff $h\left(S_{i}\right)$.
\end{itemize}
As we mentioned in the previous paragraph, we stack a few pricing
tasks together so that the agent gets a reward from the portfolio
that recently terminates. For a $T$-term option, we combine $T$
tasks together where the agent is required to learn to price not only
the $T$-term one but all $t$-term ones for $t<T$. We briefly illustrate
the stucture of RLOP in Fig. \ref{fig:rlop-illustration}.

\begin{figure}[]
\begin{centering}
\includegraphics[width=0.6\linewidth]{backup/rlop-illustration}
\par\end{centering}
\caption{A brief illustration on the RLOP environment.}

\label{fig:rlop-illustration}
\end{figure}


\subsection{Specification}

In this section, we describe how to set-up the RLOP environment and
the necessary numerical procedures.

\subsubsection{Environment}

As in the QLBS model, the environment samples a family of asset price
trajectories and maintains the portfolio value process. The asset
value is assumed to obey the geometrical brownian motion, as in Sec.
\ref{subsec:Environment}. We also allows the environment parameters
to change at the beginning of each episode according to a Poisson
process with intensity $\varUpsilon$.

\subsubsection{Agent Parametrization and Learning Algorithm}

The agent, i.e. the hedging strategy is parametrized by a Resnet-motivated
structure, same in Sec. \ref{subsec:Agent-Parametrization}. The REINFORCE
algorithm is adopted to train the agent. One caveat is that there
is no need to use a baseline as in the QLBS model, since in this case
the reward function is for penalty rather than learning the optimal
price.

\subsection{Experiments and Results}

We include a variety of experiments, including a demonstration that
shows the agent learns over time, a comparison with the Black-Scholes
baseline model, and a another one on the effect of trading costs.

\subsubsection{Demonstration on learning\label{subsec:Demonstration-on-learning-rlop}}

We use a parameter setting similar to Sec. \ref{subsec:Demonstration-on-learning}
where $r=0.01,\sigma=0.1,\mu=0,T=5,K=1,S_{0}=1,\Delta t=1$. The agent
is trained under no transaction cost with a fixed initial asset value
and the training curve is flat after around 8000 steps, shown in Fig.
\ref{fig:rlop-exp1-learns}. Similar to the QLBS setting, we introduce
adjustment to the parameters (i.e. the initial asset value in the
experiments to follow) by a Poisson process. We mark where the adjustment
occurs in the bottom half of the figure by purple triangles.

We also compare the optimal hedging strategy learnt with the BS predictions.
In general, the hedging curve should have a greater slope for a smaller
remaining time which implies that the price of the option is more
sensitive to the underlying asset price. As expected, both the learnt
position and the BS position show a greater slope when approaching
the terminal time. The one learnt by our RL agent is slightly insensitve
to the time variable.

\begin{figure}[]
\caption{Episodic return improves during training. Top figure: no adjustment
on the initial price $S_{0}$; bottom figure: the initial price is
adjusted via a Poisson process.}
\includegraphics[width=0.47\linewidth]{rlop/experiment1/learning-curve.png}

\label{fig:rlop-exp1-learns}

\end{figure}

\begin{figure}[]
\caption{Optimal hedge position learnt compared with BS predictions (Eqn. \ref{eq:bs-hedge}).}\includegraphics[width=0.47\linewidth]{rlop/experiment1/hedge_position.png}

\label{fig:rlop-exp1-hedge}
\end{figure}


\subsubsection{Effect of transaction cost\label{subsec:Effect-of-transaction-rlop}}

We proceed to measure how the transaction cost can effect the optimal
hedging strategy. We use the same setting as laid out in Sec. \ref{subsec:Effect-of-transaction}
where $\epsilon$ is the friction parameter. We plot the optimal hedging
position learnt under different $\epsilon$ in Fig. \ref{fig:rlop-exp2}.
The general trend is that a large $\epsilon$ discourages hedging
amounts, since it leads to a higher trading cost on average. A higher
friction also leads to a larger distinction between positions at different
times, which implies that in an illuquid market, we need to value
the time variable more than in a liquid one.

\begin{figure}[]
\includegraphics[width=1\linewidth]{rlop/experiment2/hedge_position.png}

\caption{Optimal hedging strategy learnt under different transaction costs.}

\label{fig:rlop-exp2}
\end{figure}